\chapter{Your own bots}
\label{ch:bots}

If you write chat bots, for example for quizzes, reminders or exercise submissions, the
admin bot can run them for you. You write a \emph{maubot plugin}, send it to the admin bot,
and after an admin's approval it runs under its own account that you invite into your
rooms. You never need server access or passwords.

\begin{note}
This chapter is only relevant if you program your own bot. The plugin framework is
\emph{maubot}; its documentation is at \url{https://docs.mau.fi/maubot/}. Ask the admins
whether professor bots are enabled on the server; otherwise the \cmd{!bot} commands answer
that they are not.
\end{note}

\section{Build the plugin}

A maubot plugin is a folder with a \file{maubot.yaml} and your Python code. The tool
\cmd{mbc} (installed with \cmd{pip install maubot}) packs it into an \file{.mbp} file:

\begin{codefile}[maubot.yaml]
maubot: 0.6.0
id: ch.unibas.prof.example.quiz
version: 1.0.0
license: MIT
modules:
  - quiz
main_class: QuizBot
config: true
extra_files:
  - base-config.yaml
\end{codefile}

\begin{codefile}[terminal]
cd my-quiz-plugin
mbc build          # creates ch.unibas.prof.example.quiz-v1.0.0.mbp
\end{codefile}

Three rules matter on this server:

\begin{description}[style=nextline,leftmargin=1.5em,font=\normalfont\bfseries]
  \item[The plugin ID must be your own]
    The \cmd{id:} in \file{maubot.yaml} identifies the plugin on the whole server. All bots
    share one plugin list, so a second plugin with the same ID would replace the first one.
    The bot therefore remembers who uploaded an ID first and refuses it for everyone else.
    Use an ID that contains your name, e.g.\ \cmd{ch.unibas.<your.name>.quiz}. Copying an
    example plugin and keeping its ID (such as \cmd{xyz.maubot.echo}) is the most common
    reason for a rejected upload.
  \item[Only packages that come with maubot]
    Plugins can't install Python packages. If \file{maubot.yaml} lists \cmd{dependencies}
    that the server doesn't have, the upload is refused, and the reply names them. Ask an
    admin if you need one.
  \item[Raise the version for every update]
    An upload identical to the running one is refused (\enquote{already runs exactly this
    file}). Change \cmd{version:} and build again.
\end{description}

\section{Request a bot}

\begin{you}
!bot create quiz-bot
\end{you}
\begin{botreply}
Now send me the plugin (.mbp file, built with \cmd{mbc build}) here, within 10 minutes (drag
it into the chat or use the paperclip). \cmd{!cancel} stops waiting.
\end{botreply}

Then send the \file{.mbp} file. The name is your choice: 1--24 lowercase letters, digits and
\cmd{-}. Each professor can have up to 3 bots (configurable).

The bot checks the file right away. If something is wrong, it says what and how to fix it,
and waits for the corrected file. You don't type the command again:

\begin{botreply}
The plugin ID \cmd{xyz.maubot.echo} is already taken by another bot on this server. All bots
share one plugin list, and uploading a plugin with the same ID would replace the other one,
so I refused it.\\
\textbf{How to fix it:} open \cmd{maubot.yaml} in your plugin folder, change the \cmd{id:} line
to an ID only you use, e.g.\ \cmd{id: ch.unibas.prof.example.echo}, run \cmd{mbc build} again
and send me the new .mbp file.\\[4pt]
Nothing was changed. You don't have to repeat the command: send the corrected plugin (.mbp
file, built with \cmd{mbc build}) here within 10 minutes, or \cmd{!cancel}.
\end{botreply}

If the file is fine, the request goes to the bot admins for approval. You get a message in
your chat with the admin bot when they decide. After approval the bot sets everything up:

\begin{botreply}
Your bot \textbf{quiz-bot} is running as @bot.prof.example.quiz-bot:matrix.dmi.unibas.ch
(plugin \cmd{ch.unibas.prof.example.quiz}).\\[4pt]
\textbf{Next steps}\\
-- Invite @bot.prof.example.quiz-bot:matrix.dmi.unibas.ch to a room; it joins by itself. \dots
\end{botreply}

The account is always \cmd{@bot.<your username>.<bot name>}, so everyone can see whose bot
it is.

\section{Manage your bots}

\begin{center}
\begin{tabularx}{\textwidth}{@{}lL@{}}
\toprule
\textbf{Command} & \textbf{What it does} \\
\midrule
\cmd{!bot list} & your bots with account, plugin and state \\
\cmd{!bot info <name>} & running or stopped, plugin version, number of rooms \\
\cmd{!bot config <name>} & sends the bot's settings (the plugin's config) as a YAML file \\
\cmd{!bot stop <name>} & stops the bot; it keeps its rooms and settings \\
\cmd{!bot start <name>} & starts it again \\
\cmd{!bot update <name>} & then send a new \file{.mbp}: new version (may need approval again) \\
\cmd{!bot delete <name>} & after \cmd{!confirm}: removes the bot and deactivates its account \\
\bottomrule
\end{tabularx}
\end{center}

\textbf{Changing the settings:} get the file with \cmd{!bot config quiz-bot}, edit it, upload
it into the chat, and reply to your uploaded file with \cmd{!bot config quiz-bot}.

\begin{warning}[What a bot can see]
A bot reads every message in the rooms it is invited to, like any member. Professor bots
run next to each other on a shared server, so don't put passwords or other secrets into a
bot's code or settings. Bot admins can stop all professor bots at once in an emergency.
\end{warning}
